% Generated draft: requires mathematical and source review.
\documentclass[11pt]{article}
\usepackage[margin=25mm]{geometry}
\usepackage{fontspec}
\setmainfont{Latin Modern Roman}
\usepackage{amsmath,amssymb,mathtools,unicode-math}
\setmathfont{Latin Modern Math}
\usepackage{enumitem,needspace,xurl,hyperref}
\setlength{\parindent}{0pt}
\setlength{\parskip}{5pt}
\setlength{\emergencystretch}{3em}
\newcommand{\N}{\mathbb{N}}
\newcommand{\Z}{\mathbb{Z}}
\newcommand{\Q}{\mathbb{Q}}
\newcommand{\R}{\mathbb{R}}
\newcommand{\C}{\mathbb{C}}
\newcommand{\F}{\mathbb{F}}
\newcommand{\A}{\mathbb{A}}
\newcommand{\PP}{\mathbb{P}}
\newcommand{\E}{\mathbb{E}}
\newcommand{\eps}{\varepsilon}
\newcommand{\dd}{\,\mathrm{d}}
\begin{document}
\section*{UnsolvedMath ID: 30002649}
\textbf{Source problem code:} catalogue entry
\subsection{Definitions and mathematical statement}
Put \(R=\Z[t^{\pm1}]\) and \(R_{\geq0}=\Z_{\geq0}[t^{\pm1}]\), the set of finite Laurent polynomials whose coefficients are nonnegative integers. The source denotes by \(S(\Sigma)\) the Kauffman bracket skein algebra of a surface \(\Sigma\) over \(R\). It does not specify the precise class of surfaces, the skein relations and parameter convention, or boundary and puncture conventions. Its background refers to a compact-surface formulation, but the statement itself only says ``a surface.'' These qualifications remain unresolved.

Let \(\mathcal L(\Sigma)\) denote the source's intended set of integer laminations. For a lamination written \(\mu=\{(a_i,C_i)\}_{i=1}^{r}\), the proposed element is
\[
 b_T(\mu)=\prod_{i=1}^{r}T_{a_i}(C_i)\in S(\Sigma),
\]
where each \(C_i\) is interpreted as a curve class in the skein algebra and polynomial evaluation uses its algebra multiplication. The source does not specify the admissible curves, their framing, the allowed weights, the treatment of peripheral or contractible curves, or the normalization of the Chebyshev polynomials \(T_j\). Thus this notation retains the source's intended conventions without selecting missing ones. In particular, no choice between classical and rescaled Chebyshev polynomials is made here.

\textbf{Question, conditional on the intended conventions.} For every surface \(\Sigma\) in the intended class, does the indexed family
\[
 \bigl(b_T(\mu)\bigr)_{\mu\in\mathcal L(\Sigma)}
\]
form an \(R\)-module basis of \(S(\Sigma)\) such that, for every \(\mu,\nu\in\mathcal L(\Sigma)\), its uniquely determined multiplication coefficients satisfy
\[
 b_T(\mu)b_T(\nu)
 =\sum_{\lambda\in\mathcal L(\Sigma)}
 c_{\mu,\nu}^{\lambda}(t)b_T(\lambda),
 \qquad
 c_{\mu,\nu}^{\lambda}(t)\in R_{\geq0},
\]
with only finitely many nonzero coefficients for each pair \((\mu,\nu)\)? Being an \(R\)-module basis means that every element of \(S(\Sigma)\) has a unique finite \(R\)-linear expansion in this family. The missing conventions must be recovered before this conditional formulation becomes a fully specified conjecture.
\subsection{Short English statement}
The question asks whether the Chebyshev-decorated integer laminations form a basis of the surface skein algebra whose multiplication coefficients are Laurent polynomials in \(t\) with nonnegative integer coefficients. The supplied record leaves the surface class, lamination conventions, skein parameter convention, and Chebyshev normalization unresolved. Its supplied status is partially solved; this draft does not independently verify that assessment.
\subsection{Sources}
\begin{itemize}
\item Local UnsolvedMath registry: permanent ID \(30002649\), source problem code \texttt{catalogue entry}, title ``Positive Bases for Surface Skein Algebras,'' in the supplied record from the Oberwolfach Reports Open Problems collection. Project registry locations are \texttt{data/unsolvedmath/problems.json} and \texttt{catalog/problems.json}. No individual proposer is identified; the supplied proposed year is 2014.
\item Supplied report citation: \emph{Low-dimensional Topology and Number Theory} (2015), Oberwolfach Reports, pp.\ 2115--2176, DOI \texttt{10.4171/owr/2014/38}. Supplied source URL and literature-source label ``original Oberwolfach report'': \url{https://doi.org/10.4171/owr/2014/38}. The background describes Oberwolfach Reports as documenting mathematical workshops and their open-problem discussions. The relation between the citation's 2015 date and the record's 2014 dating requires verification.
\item Supplied literature source: ``Mandel and Qin, bracelet bases,'' \url{https://arxiv.org/abs/2301.11101}.
\item Supplied literature assessment, checked exactly \texttt{2026-08-21}: status ``partially solved.'' The supplied assessment says that quantum bracelet positivity is proved for broad unpunctured marked-surface settings, while the full compact-surface formulation in the record is not covered in every case. This is supplied dated curation derived from the OWR curation, not an independently verified literature assessment; its time-sensitive status, cited claims, and applicability require independent verification.
\item The supplied statement-status label is \texttt{corrected\_verified}. Its accompanying verification note says that mathematical content was preserved while notation, typography, or omitted definitions were repaired for publication. This is a source-record label and qualification, not mathematical verification performed for this draft.
\end{itemize}
\end{document}
